% !TeX program = xelatex
\documentclass[UTF8,11pt,a4paper,fontset=none]{ctexart}
\usepackage[margin=24mm,headheight=15pt]{geometry}
\usepackage{fontspec}
\setmainfont{Liberation Serif}
\setsansfont{Liberation Sans}
\setmonofont{DejaVu Sans Mono}[Scale=0.82]
\setCJKmainfont{Noto Serif CJK SC}
\setCJKsansfont{Noto Sans CJK SC}
\setCJKmonofont{Noto Sans CJK SC}
\usepackage{amsmath,amssymb,mathtools,bm}
\usepackage{booktabs,longtable,array,tabularx}
\usepackage{graphicx,xcolor,tikz}
\usetikzlibrary{arrows.meta,positioning,calc}
\usepackage{enumitem,listings,fancyhdr}
\usepackage[unicode,colorlinks=true,linkcolor=blue!45!black,urlcolor=blue!45!black]{hyperref}
\usepackage{xurl}
\hypersetup{pdftitle={ProjectIcy：代码结构与耦合算法},pdfauthor={代码分析文档}}
\newcommand{\code}[1]{\nolinkurl{#1}}
\newcommand{\R}{\mathbb R}
\newcommand{\dd}{\,\mathrm d}
\newcommand{\clip}{\operatorname{clip}}
\newcommand{\diag}{\operatorname{diag}}
\newcommand{\tr}{\operatorname{tr}}
\newcommand{\rot}{\mathsf J}
\newcommand{\eps}{\varepsilon}
\newcommand{\src}[1]{\par\smallskip{\small\sffamily 代码对应：\code{#1}。}\par}
\newcolumntype{L}[1]{>{\raggedright\arraybackslash}p{#1}}
\numberwithin{equation}{section}
\setcounter{tocdepth}{2}
\setlist{nosep,leftmargin=2em}
\setlength{\parindent}{2em}
\setlength{\parskip}{3pt}
\setlength{\emergencystretch}{2em}
\renewcommand{\arraystretch}{1.18}
\pagestyle{fancy}
\fancyhf{}
\fancyhead[L]{\small ProjectIcy：代码结构与耦合算法}
\fancyhead[R]{\small\nouppercase{\leftmark}}
\fancyfoot[C]{\thepage}
\renewcommand{\sectionmark}[1]{\markboth{\thesection\quad #1}{}}
\lstset{basicstyle=\small\ttfamily,breaklines=true,columns=fullflexible,
  keepspaces=true,showstringspaces=false,frame=single,rulecolor=\color{black!20},
  backgroundcolor=\color{black!2},xleftmargin=0pt,xrightmargin=0pt}
\title{\vspace{-12mm}\textbf{ProjectIcy：代码结构与耦合算法}\\[4mm]
  \large 二维冰体下落、气液流动与融化的实现说明}
\author{}
\date{代码快照日期：2026 年 9 月 14 日}
\begin{document}
\maketitle
\vspace{-8mm}
\begin{abstract}
本文从当前 Python 实现出发，说明二维冰体在热水中下落与融化的完整计算过程。
流动部分以 D2Q9 压力—动量分布和守恒相场分布描述空气—水系统；
冰体以随体材料网格保存质量与显热，以刚体平动和转动描述运动；
热过程采用有限体积通量、显热—潜热分配以及材料到空间网格的守恒映射。
全文围绕状态变量、离散方程、边界处理、算子执行次序和守恒诊断展开，
并给出公式与函数的对应关系。所述方法以实际函数体为依据；连续方程用于解释离散结构，
不代替对具体数值修正的说明。
\end{abstract}
\noindent\textbf{阅读范围。}
分析对象为当前工作区中的 \code{iceflow2d/}、\code{tests/}、
\code{benchmarks/benchmark_solver.py} 与 \code{requirements.txt}。
本文独立编写，未读取或修改已有说明文档。
基准提交为 \code{97b8bf8dc587d6bce002c635ad8caf992c258930}，
同时计入工作区中示例默认终止时间为 $10\,\mathrm s$ 的已有改动。
精确文件摘要见同目录 \code{source_manifest.json}。

\begingroup
\small\linespread{1.0}\selectfont
\setlength{\parskip}{0pt}
\tableofcontents
\endgroup
\clearpage
\section{问题、变量与物理尺度}
\label{sec:model}

\subsection{计算对象与模型边界}
计算域为有壁矩形容器，内部包含空气、水和一块初始矩形冰。
冰体在重力及流体作用下平动、转动，接触热水后吸热融化。
冰的剩余材料始终共享一个刚体位姿；融水进入流体，不能重新冻结到冰上。
因此，当前程序求解的是单个可侵蚀刚体与气液流动、单向融化之间的耦合问题。

流体中空气与水均参与动量演化；热系统只演化水和原始冰材料。
水—气热界面绝热，空气温度在干单元中仅用作输出或温度读取的默认值。
空气比热和导热系数虽保留在配置中，当前随体热路径没有空气能量方程。
程序也没有蒸发、沸腾、辐射、冰体弹性变形或碎裂模型。
线性 Boussinesq 热浮力仅修正体力，不改变相场定义的材料密度。

\subsection{两套网格与三种相指示量}
空间网格记为 $\Omega_h$，节点索引 $(i,j)$，单元中心为
\begin{equation}
 \bm x_{ij}=\Delta x(i+\tfrac12,j+\tfrac12),
 \qquad 0\leq i<N_x,\quad 0\leq j<N_y.
\end{equation}
冰体材料网格记为 $\mathcal B_h$，索引 $a=(a_x,a_y)$。
它随冰运动，但其索引始终标识同一份初始材料。
材料网格不执行空间温度场的旋转插值更新；实际演化的是材料质量和显热。

应首先区分以下三个变量：
\begin{equation}
 \phi_{ij}\in[0,1],\qquad s_a\in[0,1],\qquad \chi_{ij}\in\{0,1\}.
\end{equation}
$\phi$ 是空气—水相场，$\phi=1$ 为水，$\phi=0$ 为空气；
$s_a$ 是某个初始冰材料单元的剩余固相比例；
$\chi$ 是流体求解器使用的尖锐冰体掩码。
三者分别对应 \code{water_phase}、\code{body_solid_fraction} 和 \code{solid_mask}。
输出变量 \code{liquid_fraction} 是由材料状态构造的展示量，不等于 $\phi$。

设壁面掩码为 $w_{ij}$，则活动流体集合与热水集合为
\begin{align}
 \mathcal F&=\{(i,j):w_{ij}=0,\ \chi_{ij}=0\},\\
 \mathcal W&=\{(i,j)\in\mathcal F:\phi_{ij}>\eps_\phi\},
 \qquad \eps_\phi=10^{-3}\quad\text{（默认）}.
\end{align}
相场在被冰覆盖的节点中可以保留历史值，作为再次暴露时的后备状态；
这些值不计入活动流体的水体积。

\begin{figure}[htbp]
\centering
\begin{tikzpicture}[>=Stealth,font=\small,
 box/.style={draw=black!60,rounded corners,align=center,text width=5.3cm,minimum height=14mm},
 label/.style={font=\scriptsize,fill=white,inner sep=1pt}]
 \node[box,fill=blue!5] (flow) {空间网格：气液流动\\$f_q,h_q,\phi,\bm u,p$};
 \node[box,fill=orange!6,right=27mm of flow] (rigid) {刚体位姿与质量性质\\$\bm O,\theta,\bm U,\omega,M,I,\bar{\bm X}$};
 \node[box,fill=blue!5,below=22mm of flow] (water) {空间网格：水的热状态\\$V_{ij},E_{ij}\ \longrightarrow\ T_{ij}$};
 \node[box,fill=orange!6,below=22mm of rigid] (body) {材料网格：冰的热状态\\$m_a,S_a\ \longrightarrow\ s_a,T_a$};
 \draw[->] ($(flow.east)+(0,3mm)$) -- node[label,above] {动量交换} ($(rigid.west)+(0,3mm)$);
 \draw[->] ($(rigid.west)-(0,3mm)$) -- node[label,below] {几何、壁速} ($(flow.east)-(0,3mm)$);
 \draw[->] (flow) -- node[label,left] {速度、湿区} (water);
 \draw[->] ($(water.east)+(0,3mm)$) -- node[label,above] {界面热量} ($(body.west)+(0,3mm)$);
 \draw[->] ($(body.west)-(0,3mm)$) -- node[label,below] {融水、显热} ($(water.east)-(0,3mm)$);
 \draw[->] ($(rigid.south)-(5mm,0)$) -- node[label,left] {位姿} ($(body.north)-(5mm,0)$);
 \draw[->] ($(body.north)+(5mm,0)$) -- node[label,right] {质量性质} ($(rigid.south)+(5mm,0)$);
\end{tikzpicture}
\caption{状态之间的主要依赖。融化还通过水体积投影和动量修正反馈到流动分布函数。}
\end{figure}

\subsection{二维量的单位}
所有积分量按单位出平面深度定义。令 $A=\Delta x^2$，则
\begin{equation}
 [V]=\mathrm{m^2},\quad [m]=\mathrm{kg/m},\quad
 [E]=\mathrm{J/m},\quad [I]=\mathrm{kg\,m}.
\end{equation}
例如材料单元初始质量为 $m_0=\rho_i A$。
字段名 \code{body_inertia_kg_m} 表示二维转动惯量单位 $\mathrm{kg\,m}$，
不是三维惯量的 $\mathrm{kg\,m^2}$。

热系统始终使用物理单位；流动和刚体积分主要使用格子单位。
以下用上标 $\ell$ 标识格子量。取参考密度 $\rho_0=\rho_w$，
固定参考格子速度 $u_{\mathrm{ref}}^\ell=0.1$，给定物理参考速度 $U_{\mathrm{ref}}$，则
\begin{equation}
 \Delta t=\frac{0.1\Delta x}{U_{\mathrm{ref}}},\qquad
 c_u=\frac{\Delta x}{\Delta t},\qquad
 \bm u=c_u\bm u^\ell.
 \label{eq:scales}
\end{equation}
其余换算为
\begin{align}
 \rho^\ell&=\rho/\rho_0,&
 \nu^\ell&=\nu\Delta t/\Delta x^2,&
 \bm g^\ell&=\bm g\Delta t^2/\Delta x,\\
 \sigma^\ell&=\frac{\sigma\Delta t^2}{\rho_0\Delta x^3},&
 p&=\rho_0 c_u^2 p^\ell,&
 \omega&=\omega^\ell/\Delta t,\\
 M&=\rho_0\Delta x^2 M^\ell,&
 I&=\rho_0\Delta x^4 I^\ell.&
\end{align}
格子总动量与角动量分别乘以
$\rho_0\Delta x^3/\Delta t$ 和 $\rho_0\Delta x^4/\Delta t$ 得到物理二维量。
参考速度用于确定单位映射，不是实际流速的截断值。

\src{config.py:LatticeScales; simulator.py:IceFlow2D.__init__}

\subsection{默认算例的数量级}
\label{sec:defaults}
命令行示例默认 $N_x\times N_y=100\times200$，物理域为
$0.025\times0.050\,\mathrm{m^2}$，故
\begin{equation}
 \Delta x=2.5\times10^{-4}\,\mathrm m,\quad
 \Delta t=6.25\times10^{-6}\,\mathrm s,\quad
 c_u=40\,\mathrm{m/s}.
\end{equation}
水面位于 $y=0.030\,\mathrm m$，冰边长 $0.008\,\mathrm m$，占 $32\times32$ 个材料单元；
冰初始转角为 $5^\circ$，最低角与水面的间距为 $0.0015\,\mathrm m$。
壁厚为 3 个格子，水池的活动面积为
\begin{equation}
 A_{w,0}=(100-6)(120-3)\Delta x^2=6.87375\times10^{-4}\,\mathrm{m^2}.
\end{equation}
水和冰的初始温度分别为 $90^\circ\mathrm C$、$0^\circ\mathrm C$，所有容器壁绝热。
流动密度比 $\rho_w/\rho_a=800$；冰水密度比 $\rho_i/\rho_w=0.917$。
默认物性如下，其中空气热物性保留在配置中，未进入当前热方程：
\begin{center}
\begin{tabular}{lrrrr}
\toprule
材料 & $\rho\ (\mathrm{kg/m^3})$ & $\nu\ (\mathrm{m^2/s})$
 & $c_p\ (\mathrm{J/(kg\,K)})$ & $k\ (\mathrm{W/(m\,K)})$\\
\midrule
水 & 1000 & $10^{-6}$ & 4186 & 0.60\\
冰 & 917 & --- & 2100 & 2.20\\
空气 & 1.25 & $1.5\times10^{-5}$ & 1005 & 0.026\\
\bottomrule
\end{tabular}
\end{center}
此外 $L=334000\,\mathrm{J/kg}$，$\sigma=0.072\,\mathrm{N/m}$，
$|\bm g|=9.8\,\mathrm{m/s^2}$，$\alpha_T=2.1\times10^{-4}\,\mathrm{K^{-1}}$；
默认格子界面参数为 $W=5$、$M_\phi=0.1$，Smagorinsky 系数为 $C_s=0.1$。
默认黏性对应
\begin{equation}
 \nu_w^\ell=10^{-4},\quad\nu_a^\ell=1.5\times10^{-3},\quad
 \tau_w=0.5003,\quad\tau_a=0.5045.
\end{equation}
实际剪切松弛还包含第 \ref{sec:collision} 节的数值稳定化。
\code{IceFlowConfig()} 的基础默认网格为 $200\times400$，并非命令行算例默认网格；
不能将这两组默认值混用。

\section{气液流动：双分布函数格子 Boltzmann 方法}
\label{sec:flow}
本节除特别注明外均使用格子单位，即 $\Delta x^\ell=\Delta t^\ell=1$，省略上标 $\ell$。
LBM 的基本操作是：在每个节点上调整离散方向分布，即碰撞；
随后沿各离散速度将分布传递到相邻节点，即迁移。
宏观量通过分布的低阶矩恢复。

\subsection{连续模型的解释}
相场对应的守恒 Allen--Cahn 模型可写为
\begin{equation}
 \partial_t\phi+\nabla\cdot(\phi\bm u)
 =\nabla\cdot\left[M_\phi\left(\nabla\phi-\Gamma\bm n\right)\right],
 \quad \Gamma=\frac{4\phi(1-\phi)}{W},\quad
 \bm n=\frac{\nabla\phi}{|\nabla\phi|+\eps}.
 \label{eq:ac}
\end{equation}
$W$ 是以格子数计的界面宽度，$M_\phi$ 为格子迁移率。
扩散项平滑界面，压缩项维持有限界面厚度。
该方程描述水—气界面；冰的融化源通过后续全局体积约束处理。

流动部分可用变密度、低马赫数的动量方程解释：
\begin{align}
 \rho(\phi)&=\rho_a+(\rho_w-\rho_a)\clip(\phi,0,1),\\
 \nabla\cdot\bm u&\simeq0,\\
 \partial_t(\rho\bm u)+\nabla\cdot(\rho\bm u\bm u)
 &\simeq-\nabla p_d+\nabla\cdot[\rho\nu(\nabla\bm u+\nabla\bm u^{\mathsf T})]
   +\bm F.
 \label{eq:continuum-momentum}
\end{align}
这里 $p_d=p-p_H$ 是相对冻结静水参考场的压力扰动。
式 \eqref{eq:continuum-momentum} 用于说明模型结构；实际推进还包含有限分辨率的相场通量、
矩源项、选择性松弛和投影，不能据此宣称离散系统严格满足连续不可压条件。

表面张力采用化学势力形式。定义自由能密度与化学势
\begin{align}
 \psi(\phi,\nabla\phi)&=\beta\phi^2(1-\phi)^2+\frac{\kappa}{2}|\nabla\phi|^2,\\
 \mu&=4\beta\phi(\phi-1)(\phi-\tfrac12)-\kappa\nabla^2\phi,\qquad
 \beta=\frac{12\sigma}{W},\quad\kappa=\frac32\sigma W.
\end{align}
总力密度为
\begin{equation}
 \bm F=(\rho-\rho_H)\bm g
 -b(\phi)\rho_w\alpha_T(T_w-T_{\rm ref})\bm g+\mu\nabla\phi,
 \quad b(\phi)=\clip(2\phi-1,0,1).
 \label{eq:force}
\end{equation}
门函数 $b$ 限制热浮力作用于水侧，避免干单元的默认空气温度驱动弥散界面。
存储的 \code{fluid_acceleration_lattice} 为 $\bm a=\bm F/\rho$。

\subsection{D2Q9 速度与平衡分布}
\begin{equation}
 \begin{array}{c|rrrrrrrrr}
 q&0&1&2&3&4&5&6&7&8\\\hline
 c_{qx}&0&1&0&-1&0&1&-1&-1&1\\
 c_{qy}&0&0&1&0&-1&1&1&-1&-1
 \end{array}
\end{equation}
权重为 $w_0=4/9$，$w_{1\ldots4}=1/9$，$w_{5\ldots8}=1/36$，
格子声速平方 $c_s^2=1/3$。反向索引满足 $\bm c_{\bar q}=-\bm c_q$。

程序中的 $f_q$ 为压力—动量分布，$h_q$ 为相场分布。
令
\begin{equation}
 P_q(\bm u)=3\bm c_q\cdot\bm u+\frac92(\bm c_q\cdot\bm u)^2-\frac32|\bm u|^2.
\end{equation}
两种平衡态为
\begin{align}
 f_q^{\rm eq}(p_d,\rho,\bm u)
 &=3(w_q-\delta_{q0})p_d+w_q\rho P_q(\bm u),\label{eq:feq}\\
 h_q^{\rm eq}(\phi,\bm u)&=w_q\phi[1+P_q(\bm u)].\label{eq:heq}
\end{align}
由 D2Q9 各向同性关系可得
\begin{align}
 \sum_q f_q^{\rm eq}&=0,&
 \sum_q\bm c_qf_q^{\rm eq}&=\rho\bm u,&
 \sum_q\bm c_q\bm c_qf_q^{\rm eq}&=p_d\bm I+\rho\bm u\bm u,\\
 \sum_qh_q^{\rm eq}&=\phi,&
 \sum_q\bm c_qh_q^{\rm eq}&=\phi\bm u.&
\end{align}
因此，$\sum f_q$ 不表示材料密度，压力也不能由 $p=c_s^2\rho$ 恢复。
压力与相场密度解耦是此处处理大密度比的关键。

\src{lattice.py:lattice_direction, lattice_weight, opposite_direction_index, momentum_equilibrium, phase_equilibrium}

\subsection{梯度、拉普拉斯与近壁延拓}
记沿方向 $q$ 一格和两格的邻值为 $\phi_q^{(1)},\phi_q^{(2)}$，本地值为 $\phi$。
代码使用两种梯度：
\begin{align}
 G_1\phi&=3\sum_qw_q\bm c_q(\phi_q^{(1)}-\phi),\\
 G\phi&=(1-\eta)G_1\phi+\frac{3\eta}{2}
   \sum_qw_q\bm c_q(4\phi_q^{(1)}-\phi_q^{(2)}-3\phi),
   \quad\eta=\frac13,\\
 L\phi&=6\sum_qw_q(\phi_q^{(1)}-\phi).
\end{align}
表面张力与相场法向使用 $G$，密度梯度源项和压力恢复使用 $G_1$，化学势使用 $L$。
化学势多项式中的相场取 $\clip(\phi,0,1)$；邻域差分读取当前存储的相场。
若第一邻点不是活动流体，则返回本地值；
若第一邻点有效而第二邻点无效，两格采样仍返回本地值。
这是 \code{_phase_neighbor} 的具体延拓规则，没有另行求解接触角边界条件。

\subsection{动量的原点矩 MRT 碰撞}
\label{sec:collision}
定义原点矩
\begin{equation}
 m_{ab}=\sum_qc_{qx}^{a}c_{qy}^{b}f_q,
 \qquad (a,b)\in\{0,1,2\}^2.
\end{equation}
尽管辅助函数名含有 \code{central_moment}，
\code{_collide_momentum} 实际传入零平移速度，以原点矩重构分布。
它与下一节的真正中心矩相场碰撞不同。

以存储动量速度 $\bm v$ 和加速度 $\bm a$ 构造受力速度
\begin{equation}
 \bm v=\frac{\sum_q\bm c_q f_q}{\rho},\qquad
 \bm u=\bm v+\frac12\bm a.
 \label{eq:physical-u}
\end{equation}
本节平衡态与源项均在该 $\bm u$ 下计算。分布源项为
\begin{equation}
 R_q=3w_q\left[\bm c_q\cdot\bm F+
 (\rho_w-\rho_a)(\bm c_q\cdot\bm u)(\bm c_q\cdot G_1\phi)\right].
 \label{eq:momentum-source}
\end{equation}
令 $m^{\rm eq}$、$R_{ab}$ 分别为平衡态和 $R_q$ 的同阶矩。
对矩向量
$\bm m=(m_{00},m_{10},m_{01},m_{20}+m_{02},m_{20}-m_{02},m_{11},m_{21},m_{12},m_{22})^{\mathsf T}$，
碰撞统一写为
\begin{equation}
 \bm m^+=\bm m-\mathsf S(\bm m-\bm m^{\rm eq})+
 (\mathsf I-\tfrac12\mathsf S)\bm R,
 \label{eq:mrt}
\end{equation}
其中
\begin{equation}
 \mathsf S=\diag(\omega_p,\omega_p,\omega_p,1,\omega_s,\omega_s,1,1,1).
\end{equation}
程序只需读取实际受非单位松弛率影响的原始矩；松弛率为 1 的矩直接取
$m^{\rm eq}+R/2$。

物理黏性与 Smagorinsky 项确定
\begin{align}
 \nu_{\rm loc}&=\bar\phi\nu_w+(1-\bar\phi)\nu_a+\nu_{\rm SGS},\qquad
 \bar\phi=\clip(\phi,0,1),\\
 \nu_{\rm SGS}&=C_s^2\sqrt{2\left[(\partial_xv_x)^2+(\partial_yv_y)^2+
 \tfrac12(\partial_yv_x+\partial_xv_y)^2\right]},\\
 \tau_p&=\frac12+3\nu_{\rm loc},\qquad\omega_p=\tau_p^{-1}.
\end{align}
这里速度梯度使用 D2Q9 一格差分；冰内邻点取刚体壁速，固定壁外取零。
额外剪切松弛下界为
\begin{equation}
 \tau_s=\max\left\{\tau_p,\frac12+(\tau_A-\tfrac12)
 \max(1-\bar\phi,\chi_{\rm nbr})\right\},\qquad \omega_s=\tau_s^{-1}.
\end{equation}
$\chi_{\rm nbr}=1$ 表示八邻域内存在尖锐冰节点，默认 $\tau_A=0.8$。
该下界只作用于偏应力矩 $m_{20}-m_{02}$ 和 $m_{11}$；
它是数值稳定化，改变局部剪切耗散，应与材料黏性区分。
设 \code{air_interface_relaxation_time=None} 时，这一额外下界关闭。

\src{simulator.py:_fluid_force_viscosity_and_gradient, _collide_momentum}

\subsection{相场的中心矩碰撞}
定义平移后的中心矩
\begin{equation}
 k_{ab}=\sum_q(c_{qx}-u_x)^a(c_{qy}-u_y)^b h_q,
 \qquad Q_q=w_q\Gamma(\bm c_q\cdot\bm n).
\end{equation}
以相同变换构造 $k^{\rm eq}_{ab}$ 和 $Q_{ab}$，并令
\begin{equation}
 \omega_\phi=(3M_\phi+\tfrac12)^{-1}.
\end{equation}
两个一阶中心矩按
\begin{equation}
 k_{ab}^+=k_{ab}-\omega_\phi(k_{ab}-k_{ab}^{\rm eq})
 +(1-\tfrac12\omega_\phi)Q_{ab},\quad (a,b)=(1,0),(0,1)
\end{equation}
更新；其余七个矩取
\begin{equation}
 k_{ab}^+=k_{ab}^{\rm eq}+\frac12Q_{ab}.
\end{equation}
特别地 $Q_{00}=0$，故零阶矩保持为 $\phi$。
重构时使用非零中心 $\bm u$；完整逆变换见附录 \ref{app:inverse}。

\src{simulator.py:_collide_phase; lattice.py:reconstruct_central_moment_population}

\subsection{迁移、宏观恢复与时间层}
对目标节点仍为流体的方向，执行推式迁移
\begin{equation}
 f_q(\bm x+\bm c_q,t+1)=f_q^+(\bm x,t),\qquad
 h_q(\bm x+\bm c_q,t+1)=h_q^+(\bm x,t).
\end{equation}
壁面链路的反射值在下一节定义。
每个方向槽有唯一写入来源，碰撞缓存使迁移不覆盖尚未使用的碰撞前状态。
迁移后先恢复
\begin{equation}
 \phi^{n+1}=\sum_qh_q^{n+1},\qquad
 \bm v^{n+1}=\frac{\sum_q\bm c_qf_q^{n+1}}{\rho(\phi^{n+1})},
\end{equation}
再恢复压力
\begin{equation}
 p^{n+1}=p_H+\frac35\left[
 \sum_{q\ne0}f_q^{n+1}+\frac12\bm u\cdot G_1\rho
 -\frac32w_0\rho|\bm u|^2\right],
 \quad G_1\rho=(\rho_w-\rho_a)G_1\phi.
 \label{eq:pressure}
\end{equation}
此处 $\bm u=\bm v^{n+1}+\bm a^n/2$，加速度来自本步碰撞前的计算，
在迁移后暂存并用于压力恢复与热对流；下一步碰撞才重新计算力。
后续几何改变、分布修正可能再次改变 $\phi$ 和 $\bm v$。

\src{simulator.py:_stream, _update_streamed_macroscopic_fields, _update_pressure}

\section{冰体几何、边界反射与刚体运动}
\label{sec:body}

\subsection{材料坐标、参考原点与质心}
定义二维旋转矩阵及垂直算子
\begin{equation}
 R(\theta)=\begin{pmatrix}\cos\theta&-\sin\theta\\\sin\theta&\cos\theta\end{pmatrix},
 \qquad \rot=\begin{pmatrix}0&-1\\1&0\end{pmatrix}.
\end{equation}
本节位置、速度、质量和惯量均采用格子单位。
材料单元中心为
\begin{equation}
 \bm X_a=(a_x+\tfrac12-N_{bx}/2,\ a_y+\tfrac12-N_{by}/2).
\end{equation}
参考原点 $\bm O$ 和转角 $\theta$ 定义材料到空间的映射；
材料质心 $\bar{\bm X}$ 与空间质心 $\bm C$ 满足
\begin{equation}
 \bm x=\bm O+R\bm X,\qquad
 \bm X=R^{\mathsf T}(\bm x-\bm O),\qquad
 \bm C=\bm O+R\bar{\bm X}.
 \label{eq:pose}
\end{equation}
非均匀融化会改变 $\bar{\bm X}$，因此 $\bm O$ 与 $\bm C$ 不能混用。
代码以 \code{body_reference_origin} 保存权威位姿，
\code{body_center} 是由式 \eqref{eq:pose} 计算的视图。
刚体在任一点的速度为
\begin{equation}
 \bm u_b(\bm x)=\bm U+\omega\rot(\bm x-\bm C).
 \label{eq:wall-velocity}
\end{equation}

\subsection{连续覆盖率与尖锐掩码}
将材料固相分数 $s_a=m_a/(\rho_i\Delta x^2)$ 双线性插值到空间单元中心，得
\begin{equation}
 C_{ij}=\clip\!\left(\sum_{a\in\mathcal N_{ij}}W_{ij,a}s_a,0,1\right),
 \quad W_{ij,a}=\max(0,1-|\xi_x-a_x|)\max(0,1-|\xi_y-a_y|).
\end{equation}
$\bm\xi$ 是材料网格的连续索引，$\mathcal N_{ij}$ 至多包含四个材料节点。
壁面处最终将 $C_{ij}$ 清零。

代码中的 \code{body_signed_distance_m} 是用于判号和切链定位的混合几何量。
在初始矩形内部，定义
\begin{equation}
 d_{ij}=\Delta x(s_*-C_{ij}),\qquad s_*=0.5\quad\text{（默认）};
 \label{eq:pseudo-sdf}
\end{equation}
矩形外部使用解析旋转矩形距离
\begin{equation}
 d_{\rm box}=\Delta x\left[|\max(\bm\delta,0)|+
 \min\{\max(\delta_x,\delta_y),0\}\right],\quad
 \bm\delta=|\bm X|-(N_{bx}/2,N_{by}/2).
\end{equation}
式 \eqref{eq:pseudo-sdf} 是覆盖率差乘以网格尺度，通常不满足 $|\nabla d|=1$；
不能将其解释为侵蚀界面的精确欧氏距离。
当刚体在力学上有效且不在壁内时，$d_{ij}\leq0$ 定义尖锐冰节点。

体积统计还需区分连续材料体积和尖锐格点体积：
\begin{equation}
 V_i^{\rm mat}=\sum_a\frac{m_a}{\rho_i},\qquad
 V_i^{\rm sharp}=\Delta x^2\sum_{ij}\chi_{ij}.
\end{equation}
前者用于质量与融化率；后者反映流动边界的离散拓扑，二者无需逐步相等。

\src{simulator.py:_body_local_coordinates, _sample_body_fraction, _signed_distance_at, _update_solid_mask}

\subsection{切链边界与动量交换}
从流体节点 $\bm x$ 沿 $\bm c_q$ 指向固体的链路称为切链。
代码使用截断后的交点比例
\begin{equation}
 \lambda_q=\clip\!\left(\frac{\max(d_f,10^{-6})}
 {\max(d_f,10^{-6})-d_s+10^{-12}},0.05,0.95\right),\qquad
 \bm x_b=\bm x+\lambda_q\bm c_q.
\end{equation}
距离 $d$ 在此以米计；$\bm x_b$ 仍用格子坐标。
截断保证交点不落在链路两端。
令 $\widehat f_q$ 为尚未施加边界处理的碰撞后分布，
以本地碰撞前的非平衡量构造反射分布
\begin{equation}
 f_{\bar q}^{\rm ref}=
 \frac{f_{\bar q}^{\rm eq}(p_d,\rho,\bm u_b)
 +[f_q-f_q^{\rm eq}(p_d,\rho,\bm u)]
 +\lambda_q\widehat f_{\bar q}}{1+\lambda_q}.
 \label{eq:cutlink}
\end{equation}
该值暂存于切链的出射槽 $q$，由迁移核写回本节点的 $\bar q$ 槽。
相场在冰面使用移动壁修正
\begin{equation}
 h_{\bar q}^{n+1}(\bm x)=h_q^+(\bm x)
 -6w_q\phi(\bm c_q\cdot\bm u_b).
\end{equation}
固定容器壁采用本地反弹；冰面额外采用式 \eqref{eq:cutlink}。

每条冰面切链向刚体贡献格子冲量
\begin{align}
 \Delta\bm P_q={}&(\bm c_q-\bm u_b)\widehat f_q
 -(-\bm c_q-\bm u_b)f_{\bar q}^{\rm ref}
 +6w_qp_H(\bm x_b)\bm c_q,\label{eq:impulse}\\
 \Delta L_q={}&(\bm x_b-\bm C)\times\Delta\bm P_q,
 \qquad \bm a\times\bm b=a_xb_y-a_yb_x.
\end{align}
$p_H(\bm x_b)$ 由链路两端线性插值。
先在节点内求和，再以双精度归约至 \code{hydrodynamic_impulse} 和
\code{hydrodynamic_torque}。

\subsection{冻结静水参考场}
初始化相场后，在每一行的活动流体节点上平均密度，得到 $\rho_H(j)$。
以初始自由液面 $j_s$ 为参考，构造
\begin{align}
 p_H(j_s-1)&=-\tfrac12g_y\rho_H(j_s-1),&
 p_H(j_s)&=\tfrac12g_y\rho_H(j_s),\\
 p_H(j)&=p_H(j+1)-\tfrac12g_y[\rho_H(j)+\rho_H(j+1)],&&j<j_s-1,\\
 p_H(j)&=p_H(j-1)+\tfrac12g_y[\rho_H(j-1)+\rho_H(j)],&&j>j_s.
\end{align}
参考场在后续时间步保持冻结。流动分布只承载 $p_d$，体力使用
$(\rho-\rho_H)\bm g$，切链冲量补回式 \eqref{eq:impulse} 的参考压力牵引。
三者共同组成静水基态的平衡处理；代码没有再向冰体额外叠加一个独立的
Archimedes 浮力项。

\src{simulator.py:_build_hydrostatic_reference, _initialize_hydrostatic_equilibrium, _prepare_momentum_cut_links}

\subsection{刚体积分与容器接触}
令总流体冲量为 $\bm J_h$，总转动冲量为 $K_h$，阻尼系数为 $d_U,d_\omega$。
每个格子时间步执行
\begin{align}
 \bm U^{n+1}&=d_U(\bm U^n+\bm J_h/M+\bm g),\\
 \omega^{n+1}&=d_\omega(\omega^n+K_h/I),\\
 \bm C^{n+1}&=\bm C^n+\bm U^{n+1},\qquad
 \theta^{n+1}=\theta^n+\omega^{n+1},\\
 \bm O^{n+1}&=\bm C^{n+1}-R(\theta^{n+1})\bar{\bm X}.
\end{align}
默认 $d_U=1$，$d_\omega=0.999$。
该阻尼按格子步应用，改变时间步也会改变单位物理时间内的阻尼效应。
积分后清空冲量累加器。

接触支撑取自\textbf{尚未进行壁面裁剪}的当前尖锐覆盖节点的四个单元面极值，
而非初始矩形外包框。
将质心投影到使这些支撑位于容器内部的区间。
在左壁处 $U_x\leftarrow\max(U_x,0)$，右壁处
$U_x\leftarrow\min(U_x,0)$，上下壁同理。
这是一种整体位置约束与向外平动速度截断，没有求解摩擦冲量、恢复系数或接触转矩。
位置确实被修正时才重新计算覆盖率。

\subsection{新覆盖与新暴露节点}
\code{previous_solid_mask} 保存前一次提交的尖锐拓扑。
新覆盖节点保存原 $\phi$ 和扰动压力 $p_d$，速度改为壁速。
新暴露节点从八邻域中\textbf{此前已经是流体且现在仍为流体}的节点平均
$\phi,p_d,\bm v$；无合法邻点时使用被覆盖时的后备值和当前壁速。
随后加回该空间位置的 $p_H$ 并重建 $f^{\rm eq},h^{\rm eq}$。
限定读取旧流体节点，使结果不依赖同一 GPU 核中新节点的执行次序。
该局部重填会改变总体积和动量，相关修正在第 \ref{sec:constraints} 节说明。

\src{simulator.py:_integrate_rigid_ice, _solve_contact_and_rasterize, _project_moving_body_inside_container, _refill_changed_nodes}

\section{随体热模型与显热—潜热分配}
\label{sec:thermal}
本节全部使用物理单位。ALE 指材料坐标与空间坐标并用的描述：
冰的热状态在材料网格中演化，水的热状态在固定空间网格中演化，
两者通过当前位姿和覆盖率交换热量。
在不发生相变的区域，导热与对流的连续物理背景为
\begin{align}
 \rho_ic_i(\partial_t+\bm u_b\cdot\nabla)T_{\rm ice}&=\nabla\cdot(k_i\nabla T_{\rm ice}),\\
 \rho_wc_w(\partial_t+\bm u\cdot\nabla)T_{\rm water}&=\nabla\cdot(k_w\nabla T_{\rm water}).
\end{align}
随体坐标消去冰的刚体输运；相变所需热量通过后述局部能量分配扣除。
程序实际推进广延量和离散通量，界面附近使用其专门的覆盖率近似。

\subsection{广延状态与温度恢复}
每个冰材料单元保存固体质量 $m_a$ 和显热 $S_a$；
每个空间水单元保存面积 $V_i$ 和显热 $E_i$。
以熔点 $T_m$ 为显热零点，初值为
\begin{equation}
 m_a^0=\rho_iA,\quad S_a^0=m_a^0c_i(T_{i,0}-T_m),\quad
 V_i^0=A\phi_i^0\ \ (i\in\mathcal W),\quad
 E_i^0=\rho_wc_wV_i^0(T_{w,0}-T_m).
\end{equation}
温度由广延量恢复：
\begin{equation}
 T_a=\min\left(T_m,T_m+\frac{S_a}{m_ac_i}\right),\qquad
 T_i=T_m+\frac{E_i}{\rho_wc_wV_i}.
 \label{eq:temperature}
\end{equation}
仅在分母大于代码中的小量阈值时作除法。
空冰材料单元返回 $T_m$，无水单元返回配置中的初始空气温度。
使用质量与能量而非直接推进温度，可使融化、单元容量变化和能量收支共享同一组状态。

\subsection{冰内部导热}
对材料网格四邻接面，定义从 $a$ 到 $b$ 的导热功率
\begin{equation}
 \dot Q_{a\to b}=k_i\min(s_a,s_b)(T_a-T_b),\qquad
 \Delta S_a^{\rm cond}=-\delta t\sum_{b\sim a}\dot Q_{a\to b}.
\end{equation}
二维面长 $\Delta x$ 与中心距离 $\Delta x$ 相消，因此式中无额外网格因子。
两端在同一不可变质量与显热状态上计算，满足
$\dot Q_{a\to b}=-\dot Q_{b\to a}$。
\code{_body_heat_delta} 暂存导热增量，待加入冰水界面热量后统一更新材料。

\subsection{水内部导热与容器热边界}
以 $x$ 向面为例，左、右单元为 $L,R$，定义水占据比例
$a_i=\clip(V_i/A,0,1)$。两侧均为合法热水单元时
\begin{equation}
 F_{E,f}^{\rm cond}=-k_w\min(a_L,a_R)(T_R-T_L).
\end{equation}
约定面通量沿正坐标方向为正，则
\begin{equation}
 E_{ij}^{\rm new}=E_{ij}-\delta t
 (F^x_{i+1,j}-F^x_{i,j}+F^y_{i,j+1}-F^y_{i,j}).
 \label{eq:fv}
\end{equation}
同一面通量由两侧以相反符号使用，内部热交换在总和中抵消。
水—气面通量为零；冰水交换另行计算。

绝热容器壁的热通量为零。定温壁以半格距离 $\Delta x/2$ 计算入域热流
\begin{equation}
 \dot Q_{\partial\Omega\to i}=2k_wa_i(T_b-T_i).
\end{equation}
累计边界输入 $Q_\partial$ 由每个导热子步的入域功率乘以 $\delta t$ 累加。
定温边界只通过相邻水单元起作用，当前路径没有直接冰—容器导热项。
默认示例的四壁均为绝热。

\subsection{水体积与显热的成对对流}
面物理速度由两侧受力速度的算术平均给出：
\begin{equation}
 u_f=\frac{c_u}{2}(u_L^\ell+u_R^\ell).
\end{equation}
按 $u_f$ 的符号选上风单元 $D$，并同时构造体积与显热通量
\begin{equation}
 F_{V,f}=u_f\Delta x\,\clip(V_D/A,0,1),\qquad
 F_{E,f}^{\rm adv}=F_{V,f}\frac{E_D}{V_D}.
 \label{eq:paired-flux}
\end{equation}
两种广延量分别用式 \eqref{eq:fv} 的离散散度更新。
面两侧都必须属于当前湿流体集合，单元不能通过干燥气相或固体面直接输运水。
界面移动所需的容量调整由第 \ref{sec:remap} 节的重映射承担。
同温水满足 $E_i=eV_i$，成对通量在纯对流中维持这一比例，避免容量变化造成伪温变。

\src{simulator.py:_accumulate_body_conduction, _compute_water_conduction_flux_x, _compute_water_conduction_flux_y, _compute_water_advection_flux_x, _compute_water_advection_flux_y}

\subsection{冰水界面的热交换}
对每个湿流体单元 $i$，检查本单元及四个轴向邻单元的冰覆盖率。
每个被检查位置 $r$ 再映射至至多四个材料节点，故一个水线程最多形成 20 个热请求。
采用调和导热系数
\begin{equation}
 k_{wi}=\frac{2k_wk_i}{k_w+k_i},\qquad
 A_{ir}=\min\left(a_i,\min(1,C_r/s_*)\right).
\end{equation}
对仍有质量的材料节点 $a$，令双线性权重为 $W_{ra}$，请求热量为
\begin{equation}
 q_{ira}=\delta t\,k_{wi}A_{ir}\frac{W_{ra}s_a}{C_r}
 \max(0,T_i-T_a),\qquad C_r>0.
\end{equation}
这一接触系数是基于覆盖率的离散近似，未重建几何界面的精确面长。
为保证水只交出实际可用的显热，对同一水单元所有请求统一缩放：
\begin{equation}
 q_i=\sum_{r,a}q_{ira},\qquad
 \eta_i=\begin{cases}\min(1,\max(E_i,0)/q_i),&q_i>0,\\0,&q_i=0,\end{cases}
 \qquad \widehat q_{ira}=\eta_iq_{ira}.
\end{equation}
水显热减少 $\sum_{r,a}\widehat q_{ira}$，相应材料热增量增加同量。
材料端采用原子累加，水线程只更新自身显热。
该热交换只允许水向冰传热；与单向融化假设一致。
当冰与热水尚无离散接触时，空气不会为冰供热。

\subsection{局部相变更新}
设材料单元累计得到净热量 $Q_a$，其旧状态为 $(m_a,S_a)$，潜热为 $L$。
若 $Q_a\leq0$ 且仍有固体，则 $S_a' = S_a+Q_a$、$m_a'=m_a$。
若 $Q_a>0$，先补足低于熔点的显热亏损，再支付潜热：
\begin{align}
 Q_{\rm sens}&=\min[Q_a,\max(-S_a,0)],& S_a'&=S_a+Q_{\rm sens},\\
 Q_{\rm rest}&=Q_a-Q_{\rm sens},&
 \Delta m_a&=\min(m_a,Q_{\rm rest}/L),\\
 m_a'&=m_a-\Delta m_a,&
 E_a^{\rm melt}&=\max(0,Q_{\rm rest}-L\Delta m_a).
 \label{eq:melt-rule}
\end{align}
材料完全融化后，其 $m_a',S_a'$ 置零，剩余正热量 $E_a^{\rm melt}$ 随融水进入空间网格。
因此，即使剩余质量极小，也必须由实际潜热支付才能消失。
这是显热与潜热分配实现的焓法；并不存在另外一个被时间积分的统一温度—液相率场。

\subsection{融水注入与密度转换}
\begin{equation}
 \Delta V_w=\Delta m_a/\rho_w,
 \qquad \Delta E_w=E_a^{\rm melt},\qquad
 \frac{\Delta V_w}{\Delta V_i}=\frac{\rho_i}{\rho_w}.
 \label{eq:melt-volume}
\end{equation}
潜热记入累计融化能，不重复加入水显热。
材料节点优先向已记录的界面水节点注入；若有多个供热节点，
记录其线性编码的最大值，以原子最大值操作确定一个目标。
无记录目标时，在材料位置附近 $5\times5$ 网格中寻找最近合法湿单元。
资格掩码先由单独内核冻结，防止注入过程改变后续线程的搜索集合。

仍无法分配的融水进入后备池，按合法湿单元的剩余容量分配；
无剩余容量权重时按已有水体积分配，再由容量同步恢复单元限制。
完全没有合法湿单元则报错。
后备分配保持质量与显热总量，但不代表精确的局部融水轨迹。

\src{simulator.py:_exchange_interface_heat, _apply_body_heat_and_phase_change, _freeze_melt_injection_eligibility, _inject_melt_water, _distribute_unassigned_melt}

\subsection{快慢时间尺度与稳定性}
\label{sec:multirate}
每个 LBM 步执行一次快速过程：容量重映射及水的对流。
每累计 $K$ 个 LBM 步执行一次慢过程：导热、边界热交换、融化和融水注入。
默认 $K=8$；慢过程按当前固定刚体位姿推进 $K\Delta t$，不是沿区间内全部历史位姿积分。

对流步长依据\textbf{面向外速度之和}而非单个节点速度范数。
对合法湿面，令格子面速度为 $\bm u_f^\ell$，外法向为 $\bm n_{if}$，定义
\begin{equation}
 r_i=\sum_{f\subset\partial i}\max(0,\bm u_f^\ell\cdot\bm n_{if}),
 \quad r_{\max}=\max_i r_i,
 \quad \delta t_{\rm adv}^{\max}=\frac{C_{\max}\Delta t}{r_{\max}}.
\end{equation}
无流出或关闭对流时上界为无穷。导热步长上界为
\begin{equation}
 \alpha_{\max}=\max\left(\frac{k_w}{\rho_wc_w},\frac{k_i}{\rho_ic_i}\right),
 \qquad \delta t_{\rm diff}^{\max}=Fo_{\max}\frac{\Delta x^2}{\alpha_{\max}}.
\end{equation}
给定过程区间 $\Delta T$，取
\begin{equation}
 N_{\rm sub}=\max\left(1,\left\lceil\Delta T/\delta t^{\max}-10^{-14}\right\rceil\right),
 \qquad \delta t=\Delta T/N_{\rm sub}.
\end{equation}
默认 $C_{\max}=0.5$、$Fo_{\max}=0.15$，每次过程至多 64 个子步。
配置要求 $Fo_{\max}\leq1/6$，与定温壁半格距离和角点扩散系数一致。
上述是代码实际采用的显式步长控制，不能替代整体耦合方法的收敛性验证。
慢热时钟只在慢过程结束时增加；若直接调用非 $K$ 倍数的 \code{step}，
其热时钟可落后于流动时钟。

\section{体积、能量与融化动量的守恒约束}
\label{sec:constraints}
该求解器有两种不同的水体积描述：流动相场的 $A\sum_{\mathcal F}\phi_i$，
以及热系统的 $\sum_iV_i$。
几何重填、相场输运和热融化分别修改它们，因此每一步需要明确的同步顺序。

\subsection{相场水体积的熵型投影}
\label{sec:phase-projection}
以格子面积为单位，水体积目标为
\begin{equation}
 V_{w,*}^{\ell}=V_{w,0}^{\ell}+M_{i,0}^{\ell}-M_i^{\ell}
 =V_{w,0}^{\ell}+\frac{\rho_i}{\rho_w}
 (V_{i,0}^{\ell}-V_i^{\ell}).
 \label{eq:volume-target}
\end{equation}
这里 $M_i^\ell=M_i/(\rho_wA)$，故格子质量损失直接等于新增格子水体积。
初始目标在相场预热之前由尖锐水区确定，预热不能重新定义水总量。

首先将活动流体的相场截断到 $[0,1]$，将
$\phi\leq\eps_\phi$ 置零、$\phi\geq1-\eps_\phi$ 置一。
这一步移除无法分辨的相场尾部，其体积变化仍由随后约束补偿。
可调整集合 $\mathcal A$ 只包含与局部 $\phi=1/2$ 穿越相连接的弥散界面节点。
代码先在八邻域中寻找穿越种子，再仅沿
$\eps_\phi<\phi<1-\eps_\phi$ 节点扩张。
扩张次数至少为 2，取不小于
\begin{equation}
 R_b=\left\lceil\frac W4\log\frac{1-\eps_\phi}{\eps_\phi}\right\rceil
\end{equation}
的偶数，以便双缓存扩张结束于主掩码。

忽略存储取整与尾部截断时，下列二元相对熵最小化解释投影形式：
\begin{align}
 \min_{\{z_i\}}\quad&\sum_{i\in\mathcal A}
 \left[z_i\log\frac{z_i}{\phi_i}+(1-z_i)\log\frac{1-z_i}{1-\phi_i}\right],\\
 \text{约束}\quad&\sum_{i\in\mathcal A}z_i+
 \sum_{i\in\mathcal F\setminus\mathcal A}\phi_i=V_{w,*}^{\ell}.
\end{align}
取约束项符号为负，驻点满足
\begin{equation}
 \log\frac{z_i}{1-z_i}=\log\frac{\phi_i}{1-\phi_i}+\lambda,
 \qquad z_i(\lambda)=\frac{\phi_i e^\lambda}{1-\phi_i+\phi_ie^\lambda}.
 \label{eq:logistic}
\end{equation}
实现直接使用该映射，无需通用优化器。
若界面剖面为 $\phi(d)=[1+\exp(-4d/W)]^{-1}$，
则 $\lambda$ 对应法向剖面平移 $\Delta d=W\lambda/4$。
因此代码限制
\begin{equation}
 |\lambda|\leq\lambda_{\max}=4d_{\max}/W,
 \qquad d_{\max}=0.5\quad\text{（默认格子数）}.
\end{equation}

实际质量函数先将 $z_i$ 转为 \code{f32}，再规范化两端尾部，最后用 \code{f64} 求和。
对未被端点钳制的自由节点，使用导数近似
\begin{equation}
 \mathcal V'(\lambda)=\sum_{i\in\mathcal A_{\rm free}}z_i(1-z_i).
\end{equation}
这是带取整单调映射的分段平滑导数估计。
先检查两端可达体积，再采用有区间保护的 Newton 迭代；
Newton 候选越界或导数无效时退回二分。
默认最多 64 次迭代。
最终对远离端点的自由节点按 $z_i(1-z_i)$ 权重分摊存储残差，最多补偿 8 次。
接收条件为
\begin{equation}
 \left|\sum_{\mathcal F}\phi_i-V_{w,*}^{\ell}\right|
 \leq\eps_V\max(1,|V_{w,*}^{\ell}|),\qquad \eps_V=10^{-8}.
\end{equation}
没有可调界面、超出位移上限或残差无法闭合时抛出异常。
实现不会在纯水或纯气体内部添加均匀体积源。

每次 $\phi\to\phi'$ 都同步提升分布：
\begin{align}
 h_q'&=h_q+h_q^{\rm eq}(\phi',\bm v)-h_q^{\rm eq}(\phi,\bm v),\\
 f_q'&=f_q+f_q^{\rm eq}(p_d,\rho',\bm v)-f_q^{\rm eq}(p_d,\rho,\bm v).
 \label{eq:population-lift}
\end{align}
这保留原非平衡部分并保持存储速度 $\bm v$。
另以 $h_0\leftarrow h_0+\phi'-\sum_qh_q'$ 修正零阶矩取整误差。
若纯相单元的 $\sum_q|h_q|>2$，则用受力速度重新构造其平衡相场分布，
以清除大幅相消的非水动力矩。

\src{simulator.py:_clip_water_phase_for_projection, _seed_water_projection_interface, _evaluate_water_projection, _apply_water_projection, _correct_water_volume}

\subsection{热水容量与显热的守恒重映射}
\label{sec:remap}
热系统的容量同步与上面的相场投影是两个独立算子。
设同步前总水面积与显热为 $V_\Sigma=\sum_iV_i$、$E_\Sigma=\sum_iE_i$，
以当前合法湿集合计算基础容量与完整容量
\begin{equation}
 B=A\sum_{i\in\mathcal W}\phi_i,\qquad F=A|\mathcal W|.
\end{equation}
要求 $0\leq V_\Sigma\leq F$（容差内）。目标面积为
\begin{equation}
 V_i^*=\begin{cases}
 A\phi_i V_\Sigma/B,&V_\Sigma\leq B,\\
 A\left[\phi_i+(1-\phi_i)\dfrac{V_\Sigma-B}{F-B}\right],&V_\Sigma>B,
 \end{cases}\qquad i\in\mathcal W,
 \label{eq:aperture}
\end{equation}
集合外取零；退化分母由代码的小量判断处理。
此式保证目标单元面积不超过 $A$，总面积跟随热系统真实水量。
当两套水体积已一致时，$V_i^*=A\phi_i$。

令比面积显热 $e_i=E_i/V_i$，单位为 $\mathrm{J/m^3}$。
仍有水的单元保留 $e_i$；新湿单元搜索半径 1、2、3 的第一圈有水邻域，按
\begin{equation}
 \widetilde e_i=\frac{\sum_{j\in\mathcal R_i}E_j/r_{ij}^2}
 {\sum_{j\in\mathcal R_i}V_j/r_{ij}^2}
\end{equation}
重建。完全无邻域供体时使用旧全域平均 $E_\Sigma/V_\Sigma$。
供体状态在此阶段保持冻结，被新固体覆盖但仍携带旧水状态的节点也可作供体。
将重建值限制在旧湿单元比显热范围 $[e_{\min},e_{\max}]$ 内，得到
$\widetilde E_i=V_i^*\widetilde e_i$。

局部重建一般不保持全局显热。令
\begin{align}
 D&=E_\Sigma-\sum_i\widetilde E_i,\\
 H&=\sum_i(e_{\max}V_i^*-\widetilde E_i),\qquad
 C=\sum_i(\widetilde E_i-e_{\min}V_i^*).
\end{align}
最终显热取
\begin{equation}
 E_i^*=\begin{cases}
 \widetilde E_i+\min(1,D/H)(e_{\max}V_i^*-\widetilde E_i),&D>0,\ H>0,\\
 \widetilde E_i+\max(-1,D/C)(\widetilde E_i-e_{\min}V_i^*),&D<0,\ C>0,\\
 \widetilde E_i,&\text{其余情形}.
 \end{cases}
\end{equation}
在可行精确算术下，该修正恢复总能量并避免引入新温度极值。
它包含非局部热量再分配，不能等同于严格局部的几何通量重映射。
代码同时记录守恒残差与
\begin{equation}
 A_E^{\rm remap}=\sum_{\text{历次同步}}|D|,
\end{equation}
后者对应 \code{aperture_energy_correction_abs_j_m}。
即使能量残差很小，$A_E^{\rm remap}$ 较大仍表示较强的数值热量再分配。
\code{ale_*} 与 \code{phase_aperture_*} 残差读取同一个组合重映射的状态，
并非两次独立操作。

\src{simulator.py:synchronize_water_aperture, _build_phase_aperture_targets, _apply_phase_aperture_transfer}

\subsection{融化后的质量、质心和惯量}
重新以格子质量 $\widehat m_a=m_a/(\rho_wA)$ 归约
\begin{equation}
 M=\sum_a\widehat m_a,\qquad
 \bar{\bm X}=\frac{\sum_a\widehat m_a\bm X_a}{M},\qquad
 I=\sum_a\widehat m_a\left(|\bm X_a-\bar{\bm X}|^2+\frac{s_a}{6}\right).
\end{equation}
单元自转项 $\widehat m_as_a/6$ 对应边长 $\sqrt{s_a}$ 的等面积正方形近似。
若旧、新质心之差在空间坐标中为 $\Delta\bm C$，则剩余材料保留原刚体速度场：
\begin{equation}
 \bm U'=\bm U+\omega\rot\Delta\bm C,
 \qquad \Delta\bm C=R(\bar{\bm X}'-\bar{\bm X}).
\end{equation}
参考原点和转角在此热更新中固定；该关系相当于无喷射反冲的材料脱离模型。

质量减小至不可分辨，或惯量低于缩放阈值时，刚体力学状态停用，
速度与角速度归零、惯量置零；剩余热质量与能量继续保留在材料网格。
因此 \code{body_active=False} 不必然表示全部质量已经融化。
双精度归约产生的微小质量变化以及非物理质量增加会被钳制到旧质量性质。

\subsection{融化线动量与角动量闭合}
在慢过程之前记录流体分布的一阶矩总和及关于固定空间原点的角动量。
冰在质量更新中失去的量按实际存储的离散状态计算：
\begin{align}
 \Delta\bm P_b&=M\bm U-M'\bm U',\\
 \Delta L_b&=\bm C\times M\bm U+I\omega-
 (\bm C'\times M'\bm U'+I'\omega').
\end{align}
融化、尖锐节点重填和式 \eqref{eq:population-lift} 都可能改变流体动量。
待水体积投影完成后，测量暂定流体变化 $\Delta\bm P_f^0,\Delta L_f^0$，
需要补足的量为
\begin{equation}
 \bm D_P=\Delta\bm P_b-\Delta\bm P_f^0,\qquad
 D_L=\Delta L_b-\Delta L_f^0.
\end{equation}
在所有合法湿流体节点上使用权重 $w_i=\phi_i$，定义
\begin{equation}
 W_s=\sum_iw_i,\qquad
 \bm x_c=\frac{\sum_iw_i\bm x_i}{W_s},\qquad
 J_s=\sum_iw_i|\bm x_i-\bm x_c|^2.
\end{equation}
补偿节点动量
\begin{equation}
 \delta\bm p_i=w_i\left[\frac{\bm D_P}{W_s}+
 \frac{D_L-\bm x_c\times\bm D_P}{J_s}\rot(\bm x_i-\bm x_c)\right]
 \label{eq:momentum-correction}
\end{equation}
满足 $\sum_i\delta\bm p_i=\bm D_P$、
$\sum_i\bm x_i\times\delta\bm p_i=D_L$。
对 D2Q9 作零阶矩为零的提升
\begin{equation}
 \delta f_{iq}=3w_q\bm c_q\cdot\delta\bm p_i,
 \qquad\sum_q\delta f_{iq}=0,\qquad
 \sum_q\bm c_q\delta f_{iq}=\delta\bm p_i.
\end{equation}
同时令 $\bm v_i\leftarrow\bm v_i+\delta\bm p_i/\rho(\phi_i)$。
完成全湿区修正后，再执行两轮两节点的力—力偶补偿以减小单精度余量。
若没有合法湿节点，或不足以支持所需角动量修正，则报错。
这是一种全局离散动量闭合，不能解释为局部融水喷射速度的解析解。

\src{simulator.py:_update_moving_body_mass_properties, _finalize_melt_momentum_coupling, _apply_melt_momentum_correction, _apply_local_melt_momentum_correction}

\subsection{总质量、参考能量与诊断含义}
权威热质量与参考能量为
\begin{align}
 M_{\rm total}&=\sum_a m_a+\rho_w\sum_iV_i,\\
 E_{\rm total}&=\sum_aS_a+\sum_iE_i+L(M_{i,0}-M_i).
 \label{eq:energy-total}
\end{align}
这里省略了初始水已经具有的常量潜热基准。
质量残差与能量残差分别为
\begin{equation}
 R_M=M_{\rm total}(t)-M_{\rm total}(0),\qquad
 R_E=E_{\rm total}(t)-E_{\rm total}(0)-Q_\partial(t).
\end{equation}
这些量只审计冰和水，不包括空气质量、机械动能、重力势能或黏性发热。
同样，融化动量残差只审计慢热过程的转移；
不能用它证明含重力、固定壁和阻尼的全系统总动量守恒。

若绝热、冰初始不高于熔点且最终完全融化，则量热学平衡温度为
\begin{equation}
 T_\infty=T_m+\frac{E_{\rm total}(0)-LM_{i,0}}
 {(M_{w,0}+M_{i,0})c_w},
 \qquad E_{\rm total}(0)\geq LM_{i,0}.
 \label{eq:calorimetry}
\end{equation}
这是给定热量足够时的平衡参考，不能预测达到平衡或完全融化所需时间。
若冰初始正好在熔点，绝热时有
$M_{\rm melt}\leq\min(M_{i,0},E_{\rm total}(0)/L)$。
低于熔点的冰还需显热加热；非均匀温度下不能把全局初始能量直接作为任意瞬态融化上界。

\section{程序组织、运行过程与结果读取}
\label{sec:program}

\subsection{文件职责与调用入口}
\begin{longtable}{L{0.39\linewidth}L{0.54\linewidth}}
\toprule 文件 & 职责\\\midrule
\endhead
\code{iceflow2d/__main__.py} & 包命令入口，调用示例的 \code{main}。\\
\code{iceflow2d/__init__.py} & 导出配置类型；通过模块 \code{__getattr__} 延迟导入求解器。\\
\code{iceflow2d/config.py} & 物性、边界、尺度与场景配置，参数验证及配置序列化。\\
\code{iceflow2d/lattice.py} & D2Q9 常量、平衡态、矩逆变换及二维叉积。\\
\code{iceflow2d/simulator.py} & 流动、热、刚体、投影与守恒诊断的 Taichi 内核及 Python 调度。\\
\code{iceflow2d/examples/coupled_falling_ice_melting_2d.py} & 物理几何转格子参数、采样调度、运行循环、CSV/NPZ/元数据和最终图。\\
\code{iceflow2d/reporting.py} & 通用快照、物理量换算、图像与 GIF 序列、进度显示。\\
\code{iceflow2d/test_*.py} & 内核状态、边界、体积约束及热耦合的配置与 CUDA 回归测试。\\
\code{tests/test_*.py} & 配置、示例调度、输出接口和进度显示等主机侧测试。\\
\code{benchmarks/benchmark_solver.py} & 预热后的步进计时及去重 Taichi 字段有效载荷统计。\\
\code{requirements.txt} & Taichi 1.7.4、NumPy、Matplotlib 和 Pillow 依赖。\\
\bottomrule
\end{longtable}
\code{examples/__init__.py} 仅标识示例包。
求解器初始化强制使用 CUDA，并关闭 CPU 回退。
纯配置、几何和大部分输出辅助代码可在不初始化 CUDA 的情况下使用。

主调用链为
\begin{equation*}
 \texttt{main}\ \longrightarrow\ \texttt{parse\_args}\ \longrightarrow\
 \texttt{create\_config}\ \longrightarrow\ \texttt{\_snapshot\_step\_targets}
 \ \longrightarrow\ \texttt{\_run\_case}.
\end{equation*}
其中 \code{_run_case} 创建 \code{IceFlow2D}，按目标步推进，并在采样边界写出结果。

\subsection{初始化的依赖顺序}
\code{IceFlow2D.__init__} 按以下依赖关系构造一致初态：
\begin{enumerate}
 \item 验证配置并初始化 CUDA；换算格子尺度，分配流动、刚体与热状态。
 \item 初始化刚体位姿、速度、质量与惯量；以完整材料质量初始化 $s_a=1$。
 \item 计算未裁剪的空间覆盖率、处理容器接触、提交覆盖率和尖锐固体掩码。
 \item 初始化尖锐水区与两套分布，记录初始活动水体积 $V_{w,0}^{\ell}$。
 \item 只执行相场碰撞与迁移进行界面预热；该过程不推进物理时钟或刚体运动。
 \item 恢复宏观量并投影水体积，使弥散界面仍对应原初始水量。
 \item 按预热后的 $\phi$ 初始化热水体积与显热，并归约材料质量性质。
 \item 构造冻结静水参考场，将流体速度初始化为零，压力初始化为 $p_H$，分布仅承载零扰动压力。
\end{enumerate}
若在相场预热前初始化水显热，热容量会与最终相场不一致；
若预热后重新定义初始水量，则会掩盖预热引入的质量变化。
这些依赖决定了构造函数的实际顺序。

\subsection{完整时间步}
下面的伪代码与 \code{IceFlow2D.step} 的调用次序一一对应。
\begin{lstlisting}[caption={一个 LBM 步的调度；pending 记录待处理慢热步数}]
collide_momentum()
collide_phase()
stream()
update_streamed_macroscopic_fields()
update_pressure()
integrate_rigid_ice()
solve_contact_and_rasterize()
update_solid_mask()
refill_changed_nodes()
advance_moving_thermal_fast(dt)
pending += 1
if pending >= K:
    advance_moving_thermal_slow(pending * dt)
    pending = 0
correct_water_volume()
finalize_melt_momentum_coupling()
synchronize_moving_water_aperture()
steps += 1
\end{lstlisting}
快速热过程先按新几何同步容量，再对流；若有多个 Courant 子步，
中间子步后再次同步容量，末子步则交给主循环后续同步。
慢过程开始前也先同步容量，再记录流体动量；随后逐导热子步执行
冰内部导热、水内部导热与壁热、冰水热交换、材料相变、融水注入和容量同步。
导热子步间只刷新连续覆盖率，尖锐 LBM 掩码在完整慢过程后统一提交。
之后更新刚体质量性质并重填因侵蚀暴露的节点，标记待闭合的融化动量。

水体积投影必须早于融化动量闭合，因为相场密度提升本身也改变分布一阶矩。
最后的容量同步使热系统与投影后的流动湿区一致。
没有慢过程的普通步，\code{_finalize_melt_momentum_coupling} 直接返回。
\code{step(0)} 不推进；负数、布尔值和非整数步数均被拒绝。

\subsection{状态存储与 GPU 执行}
\begin{longtable}{L{0.42\linewidth}L{0.17\linewidth}L{0.32\linewidth}}
\toprule 状态 & 类型 & 作用\\\midrule
\endhead
\code{momentum_populations}, \code{phase_populations} & \code{f32} & 当前 $f,h$，形状 $(N_x,N_y,9)$。\\
两套 \code{*_stream_buffer} & \code{f32} & 碰撞后分布与迁移缓存。\\
\code{water_phase}, 速度、压力、加速度 & \code{f32} & 空间宏观状态。\\
\code{solid_mask}, \code{previous_solid_mask} & \code{i8} & 当前与上次提交的冰体拓扑。\\
\code{body_reference_origin}, 转角、速度 & \code{f32} & 权威刚体运动状态。\\
\code{body_mass_lattice}, 惯量、材料质心 & \code{f64} & 融化后质量性质。\\
\code{body_solid_mass}, \code{body_sensible_energy} & \code{f64} & 材料网格上的热状态。\\
\code{water_volume_m2}, \code{water_sensible_energy} & \code{f64} & 空间网格上的热状态。\\
面通量、热增量、守恒归约 & \code{f64} & 降低大规模求和及正负相消误差。\\
\bottomrule
\end{longtable}
分布数组的逻辑索引是 $(i,j,q)$，底层按 $q$ 分块存储，以利相邻空间线程的合并读取。
相场迁移完成后，其缓存已无活跃分布用途，$q=0,1$ 两个平面被复用为体积投影掩码。
这依赖确定的算子次序，不能同时将这些平面用于另一组尚未消费的数据。

\code{_DerivedField} 是计算视图，用于壁掩码、几何距离、温度、质心和若干残差。
设备侧通过 \code{@ti.func} 即时求值，主机侧按需读取权威状态生成 NumPy 数组；
它通常没有独立设备数组，也无需刷新副本。
少数视图提供主机写入转换，例如修改质心会反算参考原点。
\code{temperature} 等完整输出场仅支持主机采样路径。

\code{@ti.kernel} 形成设备任务边界，\code{@ti.func} 内联到调用内核。
部分全局归约按 128 线程块补齐，先以 32 线程 warp 归约，再汇总局部结果。
双精度 shuffle 拆分并恢复两个 32 位字，避免将归约值降为单精度。
其他原子归约仍可能受到执行顺序影响，不能要求不同运行逐位一致。
主机读取标量与 \code{to_numpy()} 涉及设备同步或数据传输。

若以 $N=N_xN_y$、$B=N_{bx}N_{by}$ 计，主要存储复杂度为 $O(N+B)$。
每步流动计算为 $O(9N)$，热子步为 $O(N+B)$，投影还增加若干全域扫描。
性能测试需同时报告热子步、投影工作量及是否包含输出；字段有效载荷不等于总显存占用。

\subsection{配置接口与几何约束}
\begin{longtable}{L{0.39\linewidth}L{0.54\linewidth}}
\toprule 类型或参数组 & 意义及验证\\\midrule
\endhead
\code{PhaseChangeProperties} & 熔点、三相比热与导热系数、潜热；物性须为正且有限。密度统一取流动配置。\\
\code{ThermalBoundary}, \code{ThermalBoundarySet} & 四侧绝热或定温条件；绝热的附加数值必须为零。\\
\code{ThermalConfig} & 初始温度、对流开关、快慢间隔、固相阈值、Fourier/Courant 上界及热浮力参数。\\
\code{IceFlowConfig} & 物理尺度、气液物性、界面参数、几何、刚体初值、阻尼、体积投影和输出元数据。\\
\code{create_iceflow_config} & 校验覆盖参数名，拒绝未知字段，再构造配置。\\
\code{resolution}, \code{boundary_cells} & 网格必须容纳边界层及内部单元；初始旋转冰体必须位于容器内。\\
水区与重力 & 要求水平满宽水池、有空气层，气液静水参考要求竖直重力。\\
热初值 & 初始水温不低于熔点，初始冰温不高于熔点；水气界面必须绝热。\\
模式 & 仅支持 \code{body_ale}、\code{enthalpy_fv}、\code{unified}、自由刚体、静水平衡及线性水热浮力。\\
\bottomrule
\end{longtable}
示例的 \code{derive_geometry} 要求正方形网格
$L_x/N_x=L_y/N_y$，水位和冰边长对齐格子面。
初始最低冰角与水面的间距至少为 $\max(2,\lceil W\rceil)$ 格。
示例使用 $(n+1/2)/N$ 表示需转换回整数 $n$ 的尺寸比例，
避免浮点截断将预期整数尺寸减一。
水位参数是从完整域原点量起的高度，实际水深还需减去底壁厚度。

\subsection{最小运行方式}
在仓库根目录安装依赖后，可通过包入口运行短算例：
\begin{lstlisting}[language=bash]
python -m pip install -r requirements.txt
python -m iceflow2d --end-time-s 0.01 --output-interval-s 0.005 \
  --no-plot --output-dir outputs/doc_short_run
\end{lstlisting}
去掉 \code{--no-plot} 以生成三种场序列；加入 \code{--save-npz} 保存完整采样数组。
默认完整算例终止时间为 $10\,\mathrm s$，共需 $1{,}600{,}000$ 个 LBM 步，
初次运行还包含 Taichi 编译，因此短算例适合先确认执行链。
重新使用同一输出目录会覆盖历史与元数据，序列写入器也会清理该目录的旧帧。

Python 中可直接创建求解器：
\par\noindent\begin{minipage}{\linewidth}
\begin{lstlisting}[language=Python]
from iceflow2d import IceFlow2D
from iceflow2d.examples.coupled_falling_ice_melting_2d import create_config

config = create_config()
sim = IceFlow2D(config)
sim.step(config.thermal.update_interval_lbm_steps)
totals = sim.thermal.mass_energy_totals()
fields_xy = sim.sample_thermal_fields()
print(sim.physical_time_s, totals.total_energy_j_m)
\end{lstlisting}
\end{minipage}\par
配置中的 \code{show_gui} 不会在当前示例中启动交互窗口；
该例的绘图与 NPZ 保留行为由命令行参数和 \code{_run_case} 管理。

\subsection{物理时间采样与输出结构}
给定请求时刻 $t_r$，示例将其向上对齐到完整慢热周期：
\begin{equation}
 n_r=K\max\left(1,\left\lceil\frac{t_r}{K\Delta t}-10^{-12}\right\rceil\right).
\end{equation}
去除重复步号并保留 $n=0$ 的初态。
因此实际采样时刻一般略晚于请求时刻，末帧也可能略超出请求终止时刻。
默认每 $0.01\,\mathrm s$ 一帧，对应 1600 格子步，恰为完整热周期；
默认完整算例共有 1001 帧。
进度块约取 $\lceil n_{\rm end}/2000\rceil$ 步，块划分不改变数值调度。

\begin{longtable}{L{0.37\linewidth}L{0.56\linewidth}}
\toprule 输出 & 内容\\\midrule
\endhead
\code{history.csv} & 每个采样时刻的运动、融化率、体积、能量、质量与动量残差；逐行写出并刷新。\\
\code{metadata.json} & 参数、单位映射、实际采样步号、初末统计、绝热量热学参考及图像色标。\\
\code{fields.npz} & 可选；时间序列标量场、矢量场、空间坐标和诊断量。\\
\code{velocity_frames/}, \code{velocity_field.gif} & 物理水速度幅值及箭头。\\
\code{vorticity_frames/}, \code{vorticity_field.gif} & 水区涡量 $\partial_xu_y-\partial_yu_x$。\\
\code{temperature_frames/}, \code{temperature_field.gif} & 水和冰的温度。\\
\code{coupled_falling_ice_melting_2d.png} & 最后一个快照的温度及相态图。\\
\bottomrule
\end{longtable}
设备标量场形状为 $(N_x,N_y)$；\code{reporting._capture_snapshot} 转置为
$(N_y,N_x)$。NPZ 标量序列为 $(N_t,N_y,N_x)$，矢量序列多一个末轴分量维度 2。
\code{sim.sample_thermal_fields()} 本身返回未转置的 $(N_x,N_y)$ 数组。

输出的 \code{velocity_lattice} 保存动量速度 $\bm v$；
\code{physical_velocity_lattice} 保存受力速度 $\bm u$，仍需乘 $c_u$ 才是米每秒。
水速度显示掩码要求 $\phi\geq0.5$ 且非壁、非冰。
涡量仅在中心差分两侧均为合法水单元时计算，其余位置为 NaN 掩码。
色标在运行开始时固定，图中饱和颜色不表示数值速度被截断。
GIF 每帧的默认展示时长为 0.8 秒，与模拟时刻间隔不同。

\code{_StreamingGif} 逐帧编码，内存中至多保留一个解码图像；
未启用 NPZ 时无需保留全部采样数组。启用 NPZ 后，快照保留量随帧数增长。

\subsection{如何解读融化与热输出}
由材料体积计算
\begin{align}
 f_{\rm melt}&=1-V_i^{\rm mat}/V_{i,0}^{\rm mat},\\
 \ell_{\rm eq}&=\sqrt{V_i^{\rm mat}},\qquad
 d_{\rm eq}=\tfrac12(\sqrt{V_{i,0}^{\rm mat}}-\ell_{\rm eq}),\\
 \Delta V_w&=(\rho_i/\rho_w)(V_{i,0}^{\rm mat}-V_i^{\rm mat}),\\
 \Delta V_{\rm contraction}&=(1-\rho_i/\rho_w)(V_{i,0}^{\rm mat}-V_i^{\rm mat}).
\end{align}
$d_{\rm eq}$ 是等面积正方形的等效均匀融化深度，不是每个表面的局部退缩距离。

\code{sample_world_fields} 在湿单元中以
$(\rho_wLV_i+E_i)/A$ 构造展示焓，在阈值冰区映射材料显热和已融化潜热。
这一空间采样场使用最近材料索引并含有不同的常量能量基准，
不应对其直接积分以代替式 \eqref{eq:energy-total}。
权威能量取 \code{thermal.mass_energy_totals()} 或 \code{total_enthalpy_j_m()}。
干空气中的 \code{liquid_fraction} 默认也可为 1，
必须结合 \code{phase_change_material}、$\phi$ 和尖锐冰掩码解释相态。

\subsection{测试覆盖与当前核对结果}
现有测试围绕下列可检验性质展开，测试存在本身不等于全部物理工况已验证：
\begin{longtable}{L{0.38\linewidth}L{0.55\linewidth}}
\toprule 测试文件 & 核心检验\\\midrule
\endhead
\code{test_iceflow.py} & 短程状态有限性、切链反射公式、真实质量与惯量参与积分、未裁剪接触支撑及壁速度投影。\\
\code{test_volume_projection.py} & 相场有界、体积闭合、纯相不受投影污染、正常投影保留非平衡部分。\\
\code{test_thermal_coupling.py} & 热质量与能量收支、密度转换、融化动量、局部重建温度及绝热量热学约束。\\
\code{test_gpu_state.py} & 计算视图随权威状态更新、面流出 CFL 与水量正性、无效步数拒绝。\\
\code{tests/test_moving_body_thermal.py} & 模式限制、多速率控制、物性与边界验证及可选稳定性扫描。\\
\code{tests/test_coupled_falling_ice_melting_example.py} & 无 CUDA 导入路径、几何、时间调度、诊断导出和可视化接口。\\
\code{tests/test_config_validation.py}, \code{tests/test_progress.py} & 数值输入规范化、无效值拒绝及进度接口行为。\\
\bottomrule
\end{longtable}
本次执行 \texttt{python -m unittest discover -s tests -v}，共 27 项，25 项通过。
其余两项分别断言默认终止时间为 3 秒、默认终止步为 480000；
当前工作区实际值为 10 秒与 1600000，故发生断言不一致。
这是现有测试期望与当前示例参数的差异，本文按实际代码记述，未修改测试或求解器。
本次文档任务未执行完整 CUDA 物理回归或默认长时算例。

复现实验应分别观察 $R_M,R_E$、相场体积残差、容量同步残差、融化线/角动量残差，
并结合 $A_E^{\rm remap}$ 判断守恒修正对局部温度的影响。
更改网格、参考速度、界面宽度或热更新间隔时，仍需独立检验空间与时间收敛。
代码的可选全域有限性和速度扫描默认关闭，但每步热对流的面流出率归约仍执行。

\appendix
\section{D2Q9 矩逆变换}
\label{app:inverse}
在一维速度集合 $c\in\{-1,0,1\}$ 上，以 $u$ 为中心定义
$K_a=\sum_c(c-u)^a f_c$，$a=0,1,2$。
逆变换的三个系数 $B_a(c,u)$ 为
\begin{equation}
\begin{array}{c|ccc}
 c&B_0&B_1&B_2\\\hline
 0&1-u^2&-2u&-1\\
 1&u(u+1)/2&u+1/2&1/2\\
 -1&u(u-1)/2&u-1/2&1/2
\end{array}
\end{equation}
二维逆变换是张量积：
\begin{equation}
 f_q=\sum_{a=0}^2\sum_{b=0}^2
 B_a(c_{qx},u_x)B_b(c_{qy},u_y)K_{ab}.
\end{equation}
\code{inverse_central_moment_basis} 返回 $B$，
\code{reconstruct_central_moment_population} 展开九项求和。
动量路径取 $u_x=u_y=0$，还原原点矩；相场路径取受力速度，还原中心矩。
同一辅助函数支持两种矩空间，但不能据函数名将它们归为同一种碰撞。

\section{两节点力—力偶余量补偿}
全湿区修正后，若仍有线动量余量 $\bm r_P$ 和角动量余量 $r_L$，
选取湿区线性索引最小和最大的两个不同节点 $\bm x_1,\bm x_2$，记
\begin{equation}
 \bm d=\bm x_1-\bm x_2,\quad\bm x_m=\tfrac12(\bm x_1+\bm x_2),\quad
 \bm q=\frac{r_L-\bm x_m\times\bm r_P}{|\bm d|^2}\rot\bm d.
\end{equation}
令
\begin{equation}
 \delta\bm p_1=\tfrac12\bm r_P+\bm q,\qquad
 \delta\bm p_2=\tfrac12\bm r_P-\bm q.
\end{equation}
则总增量为 $\bm r_P$，总转动增量为
$\bm x_m\times\bm r_P+\bm d\times\bm q=r_L$。
代码每次重新归约实际存储结果，共执行两轮此补偿；不应将单精度存储下的残差描述为机器意义上的恒零。

\section{核心符号与代码字段索引}
\begin{longtable}{L{0.18\linewidth}L{0.39\linewidth}L{0.35\linewidth}}
\toprule 符号 & 代码字段 & 含义\\\midrule
\endhead
$\phi$ & \code{water_phase} & 水—气相场。\\
$f_q,h_q$ & \code{momentum_populations}, \code{phase_populations} & 压力—动量分布与相场分布。\\
$\bm v$ & \code{momentum_velocity_lattice} & 分布一阶矩除以材料密度。\\
$\bm a$ & \code{fluid_acceleration_lattice} & 本步缓存的流体加速度。\\
$\bm u$ & \code{physical_velocity_lattice} & 半力修正后的格子速度。\\
$p,p_H,\rho_H$ & \code{pressure_lattice}, \code{reference_pressure_by_row}, \code{reference_density_by_row} & 总压力及冻结行参考场。\\
$\chi,w,d$ & \code{solid_mask}, \code{wall_mask}, \code{body_signed_distance_m} & 尖锐冰掩码、固定壁掩码、混合几何量。\\
$m_a,S_a$ & \code{thermal.body_solid_mass}, \code{thermal.body_sensible_energy} & 材料固体质量与显热。\\
$s_a,C_{ij}$ & \code{thermal.body_solid_fraction}, \code{thermal.world_body_solid_fraction} & 材料比例及空间覆盖率。\\
$V_i,E_i$ & \code{thermal.water_volume_m2}, \code{thermal.water_sensible_energy} & 物理水面积与显热。\\
$\bm O,\theta$ & \code{body_reference_origin}, \code{body_angle} & 材料网格权威位姿。\\
$\bar{\bm X},\bm C$ & \code{body_local_center_of_mass}, \code{body_center} & 材料质心与空间质心。\\
$\bm U,\omega$ & \code{body_velocity}, \code{body_angular_velocity} & 刚体平动、角速度，格子单位。\\
$M,I$ & \code{body_mass_lattice}, \code{body_inertia_lattice} & 当前刚体格子质量与惯量。\\
$V_{w,*}^{\ell}$ & \code{water_volume_target} & 随融化更新的水体积约束。\\
$Q_\partial$ & \code{thermal.boundary_heat_input} & 累计边界入域热量。\\
$A_E^{\rm remap}$ & \code{thermal.aperture_energy_correction_abs_j_m} & 重映射全局显热修正的绝对值累积。\\
\bottomrule
\end{longtable}
文中 $M_\phi$ 专指相场迁移率，$M$ 专指刚体质量；$C_s$ 是 Smagorinsky 系数，
$C_{ij}$ 是空间冰覆盖率，$C_{\max}$ 是 Courant 上界。

\section{按执行职责定位函数}
以下路径均相对于仓库根目录。由于文档基于工作区快照，优先按函数名定位，
以免未来插入代码后行号失效。

\begingroup\small
\begin{longtable}{L{0.24\linewidth}L{0.69\linewidth}}
\toprule 阅读目标 & 函数组与前文位置\\\midrule
\endhead
尺度与验证 & \code{config.py}: \code{LatticeScales}, 各配置类的 \code{__post_init__}, \code{create_iceflow_config}；第 \ref{sec:model}、\ref{sec:program} 节。\\
初态与调度 & \code{IceFlow2D.__init__}, \code{step}, \code{_warm_start_phase}, \code{_initialize_water_volume}；第 \ref{sec:program} 节。\\
流体局部模型 & \code{_phase_neighbor}, \code{_fluid_force_viscosity_and_gradient}, \code{_collide_momentum}, \code{_collide_phase}；第 \ref{sec:flow} 节。\\
分布迁移与恢复 & \code{_stream}, \code{_stream_phase_only}, \code{_update_streamed_macroscopic_fields}, \code{_recover_initial_fluid_moments}, \code{_update_pressure}；第 \ref{sec:flow} 节。\\
几何映射 & \code{_body_local_coordinates}, \code{_nearest_body_index}, \code{_sample_body_fraction}, \code{_calculate_moving_body_fraction}, \code{_commit_moving_body_raster}, \code{_signed_distance_at}；第 \ref{sec:body} 节。\\
边界与刚体 & \code{_cut_link_fraction}, \code{_prepare_momentum_cut_links}, \code{_integrate_rigid_ice}, \code{_solve_contact_and_rasterize}, \code{_refill_changed_nodes}；第 \ref{sec:body} 节。\\
热快慢过程 & \code{_advance_moving_thermal_fast}, \code{_advance_moving_thermal_slow}, \code{MovingBodyThermal2D.advance_fast}, \code{advance_slow}；第 \ref{sec:multirate} 节。\\
显式步长 & \code{_measure_water_outflow_rate}, \code{required_advection_substeps}, \code{required_diffusion_substeps}, \code{_validate_substep_count}；第 \ref{sec:multirate} 节。\\
热通量及相变 & \code{_accumulate_body_conduction}, 四个 \code{_compute_water_*_flux_*}, \code{_update_water_advection}, \code{_update_water_conduction}, \code{_exchange_interface_heat}, \code{_apply_body_heat_and_phase_change}；第 \ref{sec:thermal} 节。\\
融水源 & \code{_freeze_melt_injection_eligibility}, \code{_inject_melt_water}, \code{_measure_melt_fallback_weights}, \code{_apply_melt_fallback}, \code{_distribute_unassigned_melt}；第 \ref{sec:thermal} 节。\\
热容量同步 & \code{synchronize_water_aperture}, \code{_measure_aperture_partial_sums}, \code{_finish_aperture_measurement}, \code{_build_phase_aperture_targets}, \code{_apply_phase_aperture_transfer}, \code{_finish_phase_aperture_transfer}；第 \ref{sec:remap} 节。\\
相场体积约束 & \code{_correct_water_volume} 及 \code{_clip_water_phase_for_projection}, \code{_seed_water_projection_interface}, \code{_dilate_water_projection_interface}, \code{_evaluate_water_projection}, \code{_apply_water_projection}, 残差权重与补偿内核；第 \ref{sec:phase-projection} 节。\\
质量性质与融化动量 & \code{_reduce_moving_body_mass_properties}, \code{_apply_moving_body_mass_properties}, \code{_snapshot_melt_momentum_coupling}, \code{_finalize_melt_momentum_coupling} 及各 \code{_reduce_*melt_momentum}、\code{_apply_*melt_momentum_correction}；第 \ref{sec:constraints} 节。\\
统计与只读视图 & \code{mass_energy_totals}, \code{total_enthalpy_j_m}, \code{sample_world_fields}, \code{sample_thermal_fields}, \code{_DerivedField}；各 \code{_read_*}、\code{*_numpy} 实现设备或主机读取；第 \ref{sec:constraints}、\ref{sec:program} 节。\\
示例运行与导出 & 示例中的 \code{build_parser}, \code{derive_geometry}, \code{create_config}, \code{_run_case}, \code{_capture_snapshot}, \code{write_history_csv}, \code{write_fields_npz}, \code{write_metadata}, \code{write_final_plot}；第 \ref{sec:program} 节。\\
通用图像与进度 & \code{reporting.py}: \code{snapshot_arrays}, \code{_physical_water_velocity}, \code{_water_vorticity_s_1}, \code{_temperature_field_for_plot}, 三个 \code{*SequenceStream}, \code{_StreamingGif}, \code{_TerminalProgress}；第 \ref{sec:program} 节。\\
性能统计 & \code{benchmarks/benchmark_solver.py}: \code{field_inventory} 对共享字段去重；\code{main} 预热 32 步，再分段计时并调用 \code{ti.sync}；计时排除输出，不代表完整算例端到端耗时。\\
\bottomrule
\end{longtable}
\endgroup

\section{文档构建与版本核对}
本文源文件为 \code{main.tex} 与同目录的七个章节文件。
使用 XeLaTeX、CTeX、Liberation 字体、Noto CJK 中文字体和 DejaVu Sans Mono 等宽字体。
在本目录执行
\begin{lstlisting}[language=bash]
bash build.sh
\end{lstlisting}
即可生成 \code{ProjectIcy_code_algorithm_zh.pdf}；中间文件保存在 \code{.build/}。
源代码文件哈希保存在 \code{source_manifest.json}，用于判断当前代码是否仍与本文快照一致。
该清单只描述代码与依赖文件，不使用已有说明文档作为依据。

\end{document}
